\documentclass[12pt,a4paper]{article}

\usepackage{fontawesome5}
\providecommand{\faPencilAlt}{\faIcon{pencil-alt}}

\usepackage{styles/couleurs}
\usepackage{styles/mise_en_page}
\usepackage{styles/boites}
\usepackage{styles/maths}
\usepackage{styles/figures}

\usepackage{enumitem}
\usepackage{array}
\usepackage{tabularx}
\usepackage{amsmath}
\usepackage{amssymb}
\usepackage{graphicx}
\usepackage{tikz}
\usepackage{xparse}

\usetikzlibrary{calc,angles,quotes,arrows.meta}

\setlist[itemize]{leftmargin=1.4em,itemsep=0.15em,topsep=0.2em}
\setlist[enumerate]{leftmargin=*,itemsep=0.22em,topsep=0.25em}

\newcommand{\lignepointillee}{\noindent\dotfill\par\vspace{0.45em}}

\newcommand{\lignesreponse}[1]{%
\begin{tcolorbox}[
enhanced,colback=white,colframe=NeutreC!40,boxrule=0.5pt,
arc=1.5mm,left=3mm,right=3mm,top=2mm,bottom=1mm,
before skip=2pt,after skip=2pt]
\foreach \i in {1,...,#1}{\lignepointillee}
\end{tcolorbox}}

\newcommand{\titreseq}[1]{%
\begin{center}
{\sffamily\large\bfseries\color{SavoirC}#1}
\end{center}}

\begin{document}

\titreseq{Activité 5 -- Calculer l'hypoténuse}

\section*{Partie A -- Défaire une multiplication}

Dans le triangle \(MNP\) rectangle en \(N\), on donne :

\[
MN=7\text{ cm}
\qquad
\widehat M=41^\circ.
\]

On cherche l'hypoténuse \(MP\).

\begin{center}
\begin{tikzpicture}[scale=0.95]
  \coordinate (M) at (0,0);
  \coordinate (N) at (5.5,0);
  \coordinate (P) at (5.5,4.8);
  \draw[line width=0.9pt] (M)--(N)--(P)--cycle;
  \draw (5.15,0)--(5.15,0.35)--(5.5,0.35);
  \draw pic[draw=ActionC,fill=ActionC!15,angle radius=9mm,"$41^\circ$"] {angle=N--M--P};
  \node[below left] at (M) {$M$};
  \node[below right] at (N) {$N$};
  \node[above right] at (P) {$P$};
  \node[below] at ($(M)!0.5!(N)$) {$7$ cm};
  \node[above left] at ($(M)!0.5!(P)$) {$?$};
\end{tikzpicture}
\end{center}

On sait que :

\[
MN=MP\times\cos(\widehat M).
\]

Complète le programme inverse :

\[
MP=MN\div\cos(\widehat M).
\]

Puis calcule \(MP\), arrondi au dixième de centimètre.

\lignesreponse{6}

\section*{Partie B -- Deux cas à ne pas confondre}

\renewcommand{\arraystretch}{1.55}
\begin{tabularx}{\textwidth}{
|>{\raggedright\arraybackslash}p{5.5cm}
|>{\centering\arraybackslash}X
|>{\centering\arraybackslash}X|}
\hline
 & \textbf{Longueur cherchée} & \textbf{Opération}\\
\hline
Hypoténuse et angle connus & côté adjacent & \\
\hline
Côté adjacent et angle connus & hypoténuse & \\
\hline
\end{tabularx}

\newpage

Complète les deux phrases.

\[
\text{adjacent}
=
\text{hypoténuse}\times\cos(\text{angle})
\]

\[
\text{hypoténuse}
=
\text{adjacent}\div\cos(\text{angle})
\]

\section*{Partie C -- Calculer}

\begin{enumerate}
  \item Dans le triangle \(RST\) rectangle en \(R\),
        \(ST\) est l'hypoténuse, \(RT=6\) cm et \(\widehat T=55^\circ\).
        Calculer \(ST\).
  \item Dans le triangle \(UVW\) rectangle en \(V\),
        \(UW\) est l'hypoténuse, \(UV=9{,}5\) m et \(\widehat U=32^\circ\).
        Calculer \(UW\).
  \item Dans le triangle \(XYZ\) rectangle en \(Y\),
        \(XZ\) est l'hypoténuse, \(XY=2{,}4\) km et \(\widehat X=76^\circ\).
        Calculer \(XZ\).
\end{enumerate}

\lignesreponse{10}

\section*{Partie D -- Vraisemblance}

Pourquoi une hypoténuse calculée doit-elle être plus grande que le côté adjacent connu ?

\lignesreponse{3}

\erreur{Ne choisis pas l'opération uniquement parce que tu veux obtenir
un nombre plus grand ou plus petit.

Pars toujours de la relation :
\[
\text{adjacent}
=
\text{hypoténuse}\times\cos(\text{angle}).
\]}

\end{document}
